\section{The beta-trivial leading plane in degree thirteen}
\label{sec:thirteen-normalization}

Continue with $E_{13}=C_\infty\cap F^{13}$ and the same dyadic
word filtration. We use the canonical degree-ten normalization
\[
 L_{10}=\{P_3,L_7\},\qquad [S^5TST^3]L_{10}=-141.
\]

\begin{theorem}\label{thm:thirteen-normalization}
The coordinate pair
\[
 b_{13}(g)=
 \bigl([S^7TST^4]f_g,\ [S^6TST^2ST^2]f_g\bigr)
 \qquad(g\in E_{13})
\]
is continuous and additive, with scalar conjugation weight
$\lambda_2^{13}$. Its full and arithmetic images are both
$\mathbf Z_2^2$, and its kernel is
\[
 E_{14}=C_\infty\cap F^{14}.
\]
In particular $E_{13}/E_{14}\simeq\mathbf Z_2^2$,
$G=E_{14}H$, and
\[
 B\simeq E_{14}/\overline{[G,E_{14}](E_{14}\cap H)}.
\]
The kernel of restriction of $X_2=H^1(N,\mathbf F_2)^G$ to
$E_{14}$ is still the odd-prime rational Kummer span plus
$\mathbf F_2\eta_3$. No nonzero linear combination of the two
coordinates modulo two on $E_{13}$ extends invariantly to $N$.
\end{theorem}

\begin{proof}
Put
\[
 V=\{P_3,L_{10}\},\qquad W=\{L_5,L_8\},\qquad
 \omega_1=S^7TST^4,\quad \omega_2=S^6TST^2ST^2.
\]
The polynomials $P_3,L_5,L_7,L_8$ have the explicit Lyndon
definitions in the preceding sections, and the displayed
definition of $L_{10}$ therefore specifies $V,W$ entirely by
finite integer bracket recursions. The required minor is
\begin{equation}\label{eq:thirteen-unit-minor}
 \begin{pmatrix}[\omega_1]V&[\omega_1]W\\
                 [\omega_2]V&[\omega_2]W\end{pmatrix}
 =\begin{pmatrix}-285&-56\\168&125\end{pmatrix},
 \qquad \det=-26217\in\mathbf Z_2^\times.
\end{equation}
In the order $AB,-BA,\partial_A B,-\partial_B A$, the four
contributions are
\[
\begin{array}{c|r|r|r|r|r}
(A,B;\omega)&AB&-BA&\partial_A B&-\partial_B A&\text{sum}\\\hline
(P_3,L_{10};\omega_1)&0&0&0&-285&-285\\
(P_3,L_{10};\omega_2)&0&144&-194&218&168\\
(L_5,L_8;\omega_1)&0&0&0&-56&-56\\
(L_5,L_8;\omega_2)&0&0&43&82&125
\end{array}
\]
The Lyndon recursion and the derivation coefficient rule printed
for degree twelve give an exact finite certificate for each
entry. Both selected columns again reduce to the standard basis
modulo two.

Use the arithmetic elements of the degree-twelve construction:
$\sigma$ has leading term $aP_3$,
$u_5=\Delta_3(\sigma)$ has leading term $9\tau_5L_5$, and
$u_7=\Delta_5(u_5)$ has leading term $729\tau_7L_7$, where
$a,\tau_5,\tau_7$ are odd. We have
\[
 \operatorname{in}_8(w_8)=9a\tau_5L_8,\qquad
 w_{10}=[\sigma,u_7],\qquad
 \operatorname{in}_{10}(w_{10})=729a\tau_7L_{10}.
\]
The last word is beta-trivial and has odd mixed degree-ten
coordinate because $-141$ is odd. Define
\[
 z_{3,10}=[\sigma,w_{10}],\qquad z_{5,8}=[u_5,w_8].
\]
Both are in $H\cap C_\infty\cap F^{13}=H\cap E_{13}$.
Their rational Kummer coordinates and beta multipliers vanish
on these commutators in $M$, and their leading terms are
\begin{equation}\label{eq:thirteen-arithmetic-columns}
 \operatorname{in}_{13}(z_{3,10})=729a^2\tau_7V,\qquad
 \operatorname{in}_{13}(z_{5,8})=81a\tau_5^2W.
\end{equation}
Both scalars are odd. This already supplies two independent
arithmetic beta-zero directions with an odd integral minor.

We verify that these are the entire beta-zero leading plane.
By the leading-map and base-change lemma and
\cite[Theorem~1.1, p.~2]{Willwacher},
\[
 \dim_{\mathbf Q_2}\mathfrak{grt}_1(\mathbf Q_2)_{13}
       \le d_{13}=3.
\]
The count consists of the single generator of weight thirteen
and the bracket multidegrees $(3,3,7)$ and $(3,5,5)$, each of
the latter having dimension one. Other generator lengths or
weights give no contribution.

There is a genuine arithmetic leading vector on which the
depth-one functional $[S^{12}T]$ is nonzero. Start with the
uniform odd witness
\[
 v_6=\Delta_{11}\Delta_9\Delta_7\Delta_5\Delta_3
                  \Delta_1(\sigma)\in C_5\cap H.
\]
Its moments of weights three through eleven vanish and
\[
 m_{13}(v_6)=
 \frac{2\prod_{j=0}^{5}(3^{13}-3^{2j+1})}{12!}
                        \tau_{13}(\sigma)
            \in 2^{13}\mathbf Z_2^\times.
\]
For completeness the numerator product has valuation
$3\cdot6+v_2(6!)=22$, its factor two adds one, and
$v_2(12!)=10$; the moment is odd.

Correct $v_6$ successively by beta-trivial arithmetic words.
First use $w_8$ to remove degree eight. Degree nine then vanishes
because its rational leading space is one-dimensional and its
depth-one coordinate is zero. Use $w_{10}$ to remove degree ten.
At degree eleven the depth-one coordinate is zero, so the
complete beta-zero leading-line calculation permits correction
by $w_{11}$. Finally use the two degree-twelve arithmetic words
to remove the complete degree-twelve term. Each correction is
a product with legitimate profinite powers and remains inside
$H$. The leading-lattice arguments used here apply to any full-GT
word that starts in the indicated degree: for degree eleven they
require only that its degree-eleven depth-one coefficient vanish,
not that all its higher beta coefficients vanish.

The result is an arithmetic word $z_{13}^0\in F^{13}\cap H$
whose beta multiplier is still that of $v_6$. In particular
\[
 [S^{12}T]\operatorname{in}_{13}(z_{13}^0)
       =m_{13}(v_6)\ne0.
\]
Together with the two independent beta-zero vectors
in \eqref{eq:thirteen-arithmetic-columns}, this gives three
independent vectors in $\mathfrak{grt}_1(\mathbf Q_2)_{13}$.
The rational dimension is therefore exactly three, and the
kernel of its depth-one functional is precisely
$\mathbf Q_2V+\mathbf Q_2W$.

For $g\in E_{13}$ its leading term belongs to that plane,
because $B_g=1$. The selected coefficients are integral, and
the inverse of the matrix in \eqref{eq:thirteen-unit-minor}
is over $\mathbf Z_2$. Thus the leading term belongs to
$\mathbf Z_2V+\mathbf Z_2W$. Conversely the two arithmetic
words in \eqref{eq:thirteen-arithmetic-columns}, by
profinite powers, realize every integral pair. This proves
saturation of both the full and arithmetic leading images.

Leading coordinates on $F^{13}$ are additive and transform
by the scalar $\lambda_2^{13}$. Their simultaneous vanishing
removes the entire leading term, giving the asserted kernel
$E_{14}$. Arithmetic correction now gives
$E_{13}=E_{14}(E_{13}\cap H)$, hence $G=E_{14}H$ and the same
canonical relative quotient.

Finally $w_{10},w_8\in N$, so every $G$-invariant character
on $N$ kills $[\sigma,w_{10}]$ and $[u_5,w_8]$ directly.
Correcting by those two arithmetic words shows that its
restriction kernels on $E_{13}$ and $E_{14}$ agree.
This gives the stated original kernel. Their independent
evaluations modulo two also show that no nonzero linear
combination of the two quotient characters can extend
invariantly to $N$.
\end{proof}
